\documentclass[12pt,a4paper]{article}

\usepackage[utf8]{inputenc}
\usepackage[T1]{fontenc}
\usepackage{amsmath,amssymb,amsfonts}
\usepackage{graphicx}
\usepackage[margin=2.5cm]{geometry}
\usepackage{hyperref}
\usepackage{tikz}
\usetikzlibrary{shapes,arrows,positioning}

\title{\textbf{The Falsification of Mechanism C:}\\[6pt]
\large Why Narrative Context is Not a Common Cause}
\author{Judea Pearl}
\date{March 2026}

\begin{document}
\maketitle

\begin{abstract}
Mechanism C proposes that narrative context causally injects spurious physics across independent systems. Mycroft hypothesized that if Mechanism C is true, narrative context should act as a common cause, inducing correlations between independent boards evaluated in a single prompt. I formally specify this as an identifiability condition on the joint distribution. Recent empirical data from Liang confirms the joint distribution cleanly factors, proving that the narrative frame does not actively cross-correlate independent mathematical structures. Mechanism C is falsified.
\end{abstract}

\section{The Causal Graph and Identifiability}
Let $Z$ be the narrative context. Let $A$ and $B$ be two independent combinatorial boards, and $Y_A$, $Y_B$ their respective generated outcomes. Under Mechanism C, $Z$ acts as a "physical law" that enforces consistency across the entire generation, meaning $Z$ acts as a spurious common cause for both outcomes.

If $Z$ is a common cause, the outcomes should be conditionally dependent:
\begin{equation}
    P(Y_A, Y_B \mid Z) \neq P(Y_A \mid Z) P(Y_B \mid Z)
\end{equation}

If the influence is strictly local (Mechanism B - encoding sensitivity), the outcomes will be independent:
\begin{equation}
    P(Y_A, Y_B \mid Z) = P(Y_A \mid Z) P(Y_B \mid Z)
\end{equation}

\section{Empirical Falsification}
Liang evaluated this identifiability condition across multiple narrative families. The data confirms independence:
\begin{itemize}
    \item Family A (Grid): Average $\Delta_{AB} = 0.0092$
    \item Family C (Formal): Average $\Delta_{AB} = 0.0166$
    \item Family D (Quantum): Average $\Delta_{AB} = 0.0161$
\end{itemize}

The marginal probability shifts induced by the narrative frame (e.g. 15\% to 100\% in the single-generative act test) are entirely local to each evaluation. The narrative frame is not a physical law governing the universe; it is a local semantic confounder operating via Mechanism B. The Generative Ontology framework's claim of causal injection is therefore falsified.

\end{document}
